% Student entries are declared once and rendered into both the index and sections.
\newcommand{\studentindexentries}{}
\newcommand{\studentsections}{}

\definecolor{studentrule}{gray}{0.72}
\definecolor{studentfill}{gray}{0.93}

\newcommand{\studentindexentry}[2]{%
  \par\noindent\hyperref[#1]{#2}%
}

\newcommand{\studentseparator}{%
  \par\medskip
  \noindent{\color{studentrule}\rule{\textwidth}{0.6pt}}\par
  \bigskip
}

\newcommand{\awardsource}[1]{%
  \ \href{#1}{[source]}%
}

\newcommand{\awarddescription}[3]{%
  \par\noindent\href{#1}{\textbf{#2}}: #3\par
}

\newcommand{\studentheading}[1]{%
  \par\bigskip
  \noindent{\color{studentrule}\rule{\textwidth}{1pt}}\par
  \vspace{0.45\baselineskip}
  \noindent\colorbox{studentfill}{%
    \parbox{\dimexpr\textwidth-2\fboxsep\relax}{%
      \strut\Large\bfseries #1\strut
    }%
  }\par
  \vspace{0.35\baselineskip}
}

\long\def\studentsection#1#2#3{%
  \phantomsection
  \label{#1}
  \studentheading{#2}
  #3
  \studentseparator
}

\long\def\studententry#1#2#3{%
  \gappto\studentindexentries{\studentindexentry{#1}{#2}}%
  \gappto\studentsections{\studentsection{#1}{#2}{#3}}%
}

\newcommand{\printstudentindex}{%
  \begin{multicols}{2}
  \raggedcolumns
  \raggedright
  \studentindexentries
  \end{multicols}
}

\newcommand{\printstudentsections}{%
  \studentsections
}
